\uuid{SgPn}
\exo7id{4456}
\titre{exo7 4456}
\auteur{quercia}
\organisation{exo7}
\datecreate{2010-03-14}
\isIndication{false}
\isCorrection{true}
\chapitre{Série numérique}
\sousChapitre{Autre}
\cpge{mpsi-mp2i.procedes-sommatoires-discrets.a, pcsi.series-numeriques.a, ptsi.series-numeriques.a, mpsi-mp2i.procedes-sommatoires-discrets.b, pcsi.series-numeriques.b, ptsi.series-numeriques.b, bcpst1.suites-reelles, mp-mpi.series-numeriques-vectorielles.b, pc.series-numeriques, psi.series-numeriques.a, pt.series-numeriques.a, bcpst2.series-reelles}
\module{Analyse}
\niveau{L1}
\difficulte{}

\contenu{
\question{
On donne une suite de réels strictement positifs $(a_n)$, décroissante
et de limite nulle. Montrer que la série de terme général
$\frac{a_n-a_{n+1}}{a_n}$ diverge.
}
\reponse{
Méthode des rectangles~:
$\sum_{k=0}^n\frac{a_k-a_{k+1}}{a_{k+1}}\ge
 \int_{t=a_{n+1}}^{a_0}\frac{d t}t \to +\infty$ lorsque $k\to\infty$.

Si $a_k\sim a_{k+1}$ la série donnée diverge donc. Sinon, elle diverge
aussi car son terme général ne tend pas vers~$0$.
}
}
